\section{An explicit positive word for the binary tensor core}\label{app:word}

This appendix gives a fixed matrix identity over a small extension field. It
will be used only as an algebraic identity: by itself it implies no
asymptotic running-time bound. Let
\[
   \F_{16}=\F_2[t]/(t^4+t+1),\qquad \eta=t^5=t^2+t.
\]
The polynomial is irreducible: it has no root in $\F_2$, and reduction
modulo the only irreducible quadratic $X^2+X+1$ gives remainder $1$.
Consequently the quotient is a field. The relation $t^4=t+1$ gives
\[
  \eta^2+\eta+1=t^4+t+1=0,\qquad \eta^3=1,
\]
so $\F_4=\F_2(\eta)$ is its subfield. Put $a=\eta^2$ and $b=\eta$;
then $ab=1$. The element $t$ is outside $\F_4$, and in particular
$t$, $t+1$, and $t^2+t+1$ are nonzero.

All matrices below act on column vectors, and products are composed
rightmost first. Write $J$ for the $3$-by-$3$ all-ones matrix and define
\[
 C_3=\eta I_3+\eta^2J
   =\begin{pmatrix}1&a&a\\a&1&a\\a&a&1\end{pmatrix},
 \qquad A=C_3\oplus[1].
\]
Since $J^2=J$ in characteristic two, direct multiplication gives
$C_3^2=\eta^2I_3+\eta J$ and $C_3^3=I_3$. Thus $A^{-1}=A^2$.
The only nonmonomial letter permitted below is the forward matrix $A$;
every inverse occurrence is replaced by two forward occurrences. Diagonal
matrices with nonzero entries and coordinate permutations are monomial.

\paragraph{Projectors and a rank-one commutator.}
Let $E_{ij}$ denote the standard matrix units and set
\[
 u=(1,a,a,0)^T,\quad v=(1,b,b,0)^T,\quad
 P=A E_{11}A^{-1}=uv^T.
\]
Indeed, the first column of $A$ is $u$, and the first row of $A^{-1}$
is $v^T$. Also $v^Tu=1+ab+ab=1$, hence $P^2=P$. With
\[
 D=\operatorname{diag}(t,1,t+1,1),\quad Q=DPD^{-1}=x_0y_0^T,
\]
where
\[
 x_0=(t,a,a(t+1),0)^T,\qquad
 y_0=(t^{-1},b,b/(t+1),0)^T,
\]
we again have $Q^2=Q$. Calculation in $\F_{16}$ gives
\[
 v^Tx_0=t+1+(t+1)=0,\qquad
 y_0^Tu=t^{-1}+1+(t+1)^{-1}=\kappa,
 \quad \kappa=\frac{t^2+t+1}{t(t+1)}\ne0.
\]
Therefore $PQ=0$ while $QP=\kappa x_0v^T$.

For idempotents $P,Q$ satisfying $PQ=0$, expansion using
$(I+cP)^{-1}=I+\frac{c}{1+c}P$ yields
\[
 [I+cP,I+dQ]:=(I+cP)(I+dQ)(I+cP)^{-1}(I+dQ)^{-1}
      =I+\frac{cd}{1+c}QP,
\]
when the displayed inverses exist. To see the cancellation, multiply the
first three factors and use $PQ=0$: the $P$ terms cancel, leaving
$I+dQ+\frac{cd}{1+c}QP$. On multiplication by $(I+dQ)^{-1}$,
the $Q$ term cancels and the extra term vanishes because
$QPQ=Q(PQ)=0$.

Define the diagonal $R=\operatorname{diag}(\eta,1,1,1)$ and let
\[
 G=ARA^{-1}=I+aP,\qquad H=DGD^{-1}=I+aQ.
\]
The formula above with $c=d=a$ has coefficient
$a^2/(1+a)=\eta^4/(1+\eta^2)=1$. Thus
\[
 T_0=GHG^{-1}H^{-1}=I+QP=I+\kappa x_0v^T.
\]
This is already a positive word: each of $G,H,G^{-1},H^{-1}$ uses
three forward $A$ letters after replacing $A^{-1}$ by $A^2$. Hence
$T_0$ uses 12 forward $A$ letters; its other letters are monomial.

\paragraph{Four successive transvections.}
Let $B=\operatorname{diag}(1,\eta,\eta,1)$. Then
\[
 x=Bx_0=(t,1,t+1,0)^T,\qquad w^T=v^TB^{-1}=(1,1,1,0),
 \quad T_1=BT_0B^{-1}=I+\kappa xw^T.
\]
Let $S_{12}$ interchange coordinates 1 and 2. It fixes the row $w^T$
and $w^TS_{12}x=1+t+(t+1)=0$. Consequently, multiplying the two
rank-one updates in $T_1S_{12}T_1S_{12}^{-1}$ gives
\[
 T_2=I+\kappa(t+1)(e_1+e_2)w^T.
\]
In this calculation the two cross terms combine as
$\kappa (x+S_{12}x)w^T$; the quadratic term is zero because
$w^Tx=0$. Now take $D_3=\operatorname{diag}(1,1,\eta,1)$.
The row difference $w^T+w^TD_3^{-1}$ is
$(1+\eta^{-1})e_3^T=\eta e_3^T$, and both row vectors pair to zero
with $e_1+e_2$. Hence
\[
 T_3=T_2D_3T_2D_3^{-1}
    =I+\kappa(t+1)\eta(e_1+e_2)e_3^T.
\]
With $D_1=\operatorname{diag}(\eta,1,1,1)$, the corresponding
rank-one update has zero cross-pairing; its column difference is
$(1+\eta)e_1=\eta^2e_1$. Therefore
\[
 T_4=T_3D_1T_3D_1^{-1}
    =I+\kappa(t+1)E_{13}.
\]
Put $\mu=\kappa(t+1)\ne0$. Diagonal conjugation removes this
coefficient:
\[
 T_5=\operatorname{diag}(\mu^{-1},1,1,1)\,T_4\,
          \operatorname{diag}(\mu,1,1,1)
       =I+E_{13}=:U_{13}(1).
\]
Here $U_{ij}(c)=I+cE_{ij}$ is a transvection. At each of the four
commutator stages the word length doubles, while the conjugating diagonal
adds two monomial letters. The forward-$A$ counts are therefore
$12,12,24,48,96,96$ for $T_0,T_1,T_2,T_3,T_4,T_5$, respectively.
The final word for $U_{13}(1)$ has 192 letters, 96 of them equal to $A$.
Conjugating by a permutation gives a positive word for any desired
$U_{ij}(1)$, with the same 96 forward-$A$ occurrences and two additional
permutation letters except for the base pair.

\paragraph{The target and its factor order.}
Write
\[
 H_2=\begin{pmatrix}1&1\\1&0\end{pmatrix}=U_{12}(1)S,
\]
where $S$ swaps the two coordinates. Let
$P_f=(1\ 3)(2\ 4)$ interchange the first-coordinate labels and
$P_s=(1\ 2)(3\ 4)$ interchange the second-coordinate labels. With coordinates
ordered $(1,1),(1,2),(2,1),(2,2)$, the exact rightmost-first identity is
\[
 H_2\otimes H_2
   =U_{13}(1)U_{24}(1)P_f\,U_{12}(1)U_{34}(1)P_s.
\]
The first two transvections are the two disjoint copies of $U_{12}$ on
the first tensor factor, and the next two are the copies on the second;
the intervening and final permutations implement the two swaps. Equivalently
the combined wire permutation is $P_fP_s=(1\ 4)(2\ 3)$, and the four
transvections are all over $\F_2$. Replace each $U_{ij}(1)$ by the
explicit conjugate of $T_5$ just constructed. The resulting word has
exactly $4\cdot96=384$ forward $A$ letters. The four $U$ words contribute
768 letters before conjugating permutations; the three non-base
conjugations add six permutation letters, and $P_f,P_s$ add two more.
Thus the full positive word has 776 letters, of which 392 are monomial.
Its product is exactly $H_2\otimes H_2$, whose entries lie in $\F_2$.

A separate bounded witness checks the chosen coordinates and finite identities; it is corroboration for this calculation, not a proof of a runtime bound. That witness is not included in this focused bundle. The word alone
does not establish a tensor-rank saving, a fast $A$-tensor algorithm, or
an asymptotic multiplication theorem. Those claims require separate
arguments for tensoring the word, routing its monomial factors, and
charging every operation in the chosen machine model.
